Formal power series extend the coefficient description of polynomials in
\wikilink{s25--b006}{the polynomial monoid ring}. Throughout, $k$ is a field
and $t$ is a formal variable. No analytic convergence is involved.

\begin{definition}\label{formal-series}
The ring $k[[t]]$ of \emph{formal power series} consists of all sequences
$(a_0,a_1,\ldots)$ in $k$, written $\sum_{n\ge0}a_nt^n$, with
\[
 (a+b)_n=a_n+b_n,\qquad (ab)_n=\sum_{i+j=n}a_ib_j.
\]
Its identity has constant coefficient $1$ and all other coefficients zero.
Finitely supported sequences form the polynomial subring $k[t]$.
\end{definition}
\begin{proof}\label{series-ring-proof}
Each coefficient sum is finite. Both parenthesizations of a triple product
have coefficient $\sum_{i+j+\ell=n}a_ib_jc_\ell$ in degree $n$.
Distributivity and the other axioms follow coefficientwise from $k$.
Finite supports remain finite under addition and convolution, whose
formulas are exactly those for polynomial coefficients.
\end{proof}

\begin{proposition}\label{series-completion}
Truncation gives a natural ring isomorphism
$k[[t]]\cong\varprojlim_{n\ge1}k[t]/(t^n)$.
\end{proposition}
\begin{proof}\label{series-completion-proof}
Send a series to $(\sum_{i<n}a_it^i+(t^n))_n$. Conversely, compatible
residue classes have unique representatives of degree less than $n$
whose coefficients agree under truncation. They determine one series.
These inverse constructions respect the coefficient operations. This
identifies explicitly the \wikilink{s21--b017}{completion example}.
\end{proof}

\begin{theorem}\label{series-units}
A series $f=\sum a_nt^n$ is a unit if and only if $a_0\ne0$.
\end{theorem}
\begin{proof}\label{series-units-proof}
If $fg=1$, constant coefficients give $a_0b_0=1$.
Conversely, set $b_0=a_0^{-1}$ and recursively define
\[
 b_n=-a_0^{-1}\sum_{i=1}^n a_ib_{n-i}\quad(n\ge1).
\]
The series $g=\sum b_nt^n$ satisfies $fg=1$ coefficientwise, hence
$gf=1$ by commutativity.
\end{proof}

\begin{definition}\label{series-order}
For nonzero $f=\sum a_nt^n$, its \emph{order} is
$\nu(f)=\min\{n:a_n\ne0\}$. Set $\nu(0)=+\infty$.
\end{definition}
\begin{proposition}\label{series-order-properties}
For nonzero $f,g$, $\nu(fg)=\nu(f)+\nu(g)$. Thus $k[[t]]$ is an
integral domain. Every nonzero series has a unique expression
$f=t^{\nu(f)}u$ with $u$ a unit.
\end{proposition}
\begin{proof}\label{series-order-properties-proof}
If $r=\nu(f)$ and $s=\nu(g)$, coefficients below $r+s$ vanish and
the coefficient of degree $r+s$ is $a_rb_s\ne0$.
Shifting the coefficients of $f$ down by $r$ gives the unique series
$u$, whose constant coefficient is nonzero.
\end{proof}

\begin{theorem}\label{series-ideals}
The ring $k[[t]]$ is local with unique maximal ideal $(t)$.
Every nonzero ideal is $(t^n)$ for a unique $n\ge0$.
Consequently $k[[t]]$ is a principal ideal domain.
\end{theorem}
\begin{proof}\label{series-ideals-proof}
Nonunits are exactly the series in $(t)$, and evaluation of the constant
coefficient identifies the quotient with $k$. Thus $(t)$ is the unique
maximal ideal. For a nonzero ideal $I$, choose $f\in I$ of minimal
order $n$. Since $f=t^nu$ with $u$ a unit, $t^n\in I$.
Every element of $I$ has order at least $n$ and is divisible by $t^n$,
so $I=(t^n)$. Its least nonzero order is $n$, proving uniqueness.
The zero ideal is principal too, and the domain property was proved above;
see \wikilink{a2s15--b001}{the definition of a PID}.
\end{proof}
\begin{example}\label{geometric-series}
The formal identity $(1-t)\sum_{n\ge0}t^n=1$ follows by cancellation
of each positive-degree coefficient. It does not assert convergence
of a numerical series at a chosen value of $t$.
\end{example}
\begin{remark}\label{series-coefficient-scope}
Over a general commutative ring $A$, the unit criterion becomes: the
constant coefficient is a unit of $A$. The same recursion proves it.
The local and PID conclusions above use the field assumption on $k$.
\end{remark}
